\documentclass[11pt,a4paper]{article}

\usepackage[utf8]{inputenc}
\usepackage[T1]{fontenc}
\usepackage{amsmath,amssymb,amsthm}
\usepackage{geometry}
\usepackage{hyperref}
\usepackage{xcolor}
\usepackage{braket}
\usepackage{algorithm}
\usepackage{algpseudocode}
\usepackage{listings}
\usepackage{booktabs}

\geometry{margin=1in}

\newtheorem{theorem}{Theorem}
\newtheorem{proposition}[theorem]{Proposition}
\newtheorem{corollary}[theorem]{Corollary}
\newtheorem{definition}{Definition}
\newtheorem{remark}{Remark}
\newtheorem{example}{Example}
\newtheorem{conjecture}{Conjecture}

\newcommand{\Rset}{\mathbb{R}}
\newcommand{\Cset}{\mathbb{C}}
\newcommand{\Nset}{\mathbb{N}}
\newcommand{\Hil}{\mathcal{H}}
\newcommand{\Dom}{\mathcal{D}}
\newcommand{\Kop}{\hat{K}}
\newcommand{\Eset}{\mathbb{E}}

\lstset{
  basicstyle=\ttfamily\small,
  breaklines=true,
  frame=single,
  captionpos=b,
}

\title{\textbf{Resolution of Singularities in Ordinary Differential Equations\\
via Koopman--Weyl Quantization and Essential Self-Adjointness}\\[6pt]
\large A Corrected and Substantially Expanded Account with\\
Algorithmic Singularity Detection and Formal Verification}
\author{}
\date{July 2026}

\begin{document}
\maketitle

\begin{abstract}
We show that the classical initial value problem $\dot{x}=x^{2}$ admits a
continuous solution for all time when the initial condition carries any
finite amount of uncertainty---unavoidable in both engineering practice and
numerical computation.  This stands in stark contrast with the singular
deterministic solution $x(t)=x(0)/(1-t\,x(0))$, which blows up at
$t=1/x(0)$.

The resolution proceeds in three stages.  First, we quantize the ODE via
the \emph{Koopman--Weyl} prescription, obtaining a symmetric operator
$H=\tfrac{1}{2}(\hat{p}\,x^{2}+x^{2}\hat{p})$ on $L^{2}(\Rset)$.
Second, we invoke \emph{Nelson's flow-completeness criterion}: the
operator $H$ is essentially self-adjoint on $C^{\infty}_{c}(\Rset)$ if and
only if the classical flow is complete (defined for all time).  For
$\dot{x}=x^{2}$ the flow is \emph{incomplete}, so $H$ is \emph{not}
essentially self-adjoint---the deficiency indices are non-zero and the
unitary time-evolution is not uniquely determined on the original sample
space.  Third, we resolve the singularity by extending the sample space:
embedding $L^{2}(\Rset)$ into $L^{2}(\Rset^{2})$ via a complexification
$z=x+iy$ with $y(0)\neq 0$ almost surely, the singularity condition
$t=1/z(0)$ is never met except on a null-measure set.

We present the complete algorithmic pipeline implemented in the
\texttt{ode\_sirk} crate---classical flow integration, blow-up time
computation via quadrature, coordinate transformations (reciprocal and
logarithmic), Weyl-symmetrized Hamiltonian construction, and essential
self-adjointness reporting---together with its formal verification in
Lean~4 (\texttt{timepiece/Singularity}).  Five validation cases are
analyzed in detail.
\end{abstract}

\tableofcontents
\newpage

%% ====================================================================
\section{Introduction}\label{sec:intro}
%% ====================================================================

Consider the scalar autonomous ODE
\begin{equation}\label{eq:ode}
  \dot{x}=x^{2},\qquad x(0)=x_{0}\in\Rset.
\end{equation}
Its unique maximal solution is
\begin{equation}\label{eq:classical-sol}
  x(t)=\frac{x_{0}}{1-t\,x_{0}},
\end{equation}
defined on the maximal interval $(-\infty,\,1/x_{0})$ when $x_{0}>0$ and
$(1/x_{0},\,+\infty)$ when $x_{0}<0$.  In either case the solution
\emph{blows up} in finite time: $\lvert x(t)\rvert\to\infty$ as
$t\to 1/x_{0}$.  There is no global ($t\in\Rset$) deterministic
solution~\cite{GORIELY2000422}.

This finite-time blow-up creates two distinct problems:
\begin{enumerate}
  \item \textbf{Undetectability.}  If a numerical integrator fails to
    converge, one cannot distinguish a genuine singularity from
    insufficient resolution.
  \item \textbf{Degeneracy.}  The blow-up permits unrelated initial
    conditions before and after the singular time, destroying uniqueness
    of the initial value problem.
\end{enumerate}

The central claim of this article is that \emph{both problems are resolved
by admitting a finite amount of uncertainty in the initial condition}.
Concretely, we replace the point $x_{0}\in\Rset$ by a probability measure
$\mu_{0}$ on $\Rset$ (equivalently, a wavefunction
$\psi_{0}\in L^{2}(\Rset)$ with $\lvert\psi_{0}\rvert^{2}=\mu_{0}$) and
evolve it under the \emph{Koopman--Weyl} unitary group generated by the
quantized Hamiltonian.  The resulting evolution is globally defined in
time, non-deterministic in the coordinate $x$, and reduces to the
classical solution in the zero-uncertainty limit.

The article is organized as follows.
Section~\ref{sec:koopman} reviews the Koopman--von Neumann formalism and
corrects several subtleties in the original treatment.
Section~\ref{sec:weyl} constructs the Weyl-symmetrized Hamiltonian.
Section~\ref{sec:nelson} states Nelson's theorem and analyzes essential
self-adjointness.
Section~\ref{sec:diagonalization} diagonalizes the Hamiltonian for
$\dot{x}=x^{2}$ and derives the explicit unitary evolution.
Section~\ref{sec:density} gives the correct density evolution formula,
including the Jacobian factor omitted in earlier treatments.
Section~\ref{sec:resolution} explains how uncertainties resolve the
singularity.
Section~\ref{sec:algorithm} presents the algorithmic singularity detector
implemented in \texttt{ode\_sirk}.
Section~\ref{sec:cov} describes coordinate transformations that resolve
singularities.
Section~\ref{sec:pipeline} gives the full analysis pipeline.
Section~\ref{sec:lean} summarizes the Lean~4 formal verification.
Section~\ref{sec:tests} presents five validation cases.
Section~\ref{sec:degenerate} discusses degenerate solutions.
Section~\ref{sec:conjecture} states generalizations and open conjectures.

%% ====================================================================
\section{Koopman--von Neumann Formalism}\label{sec:koopman}
%% ====================================================================

\subsection{Classical statistical mechanics as quantum mechanics}

Given a sample space $\Omega=\Rset$ with a probability measure $\mu$, one
may always define a wavefunction $\psi=\sqrt{d\mu/dx}$ (up to a phase).
The Koopman--von Neumann
formulation~\cite{Sudarshan1976Interaction,PhysRevResearch.2.043102,BEVANDA2021197}
recasts classical statistical mechanics as a special case of quantum
mechanics in which the algebra of observables is commutative (because the
time-evolution is deterministic).

\begin{remark}[Commutativity caveat]
For two \emph{unbounded} self-adjoint operators, commutativity of the
operators does \emph{not} imply commutativity of their spectral
measures~\cite{schmuxfcdgen2012unbounded}; one says they \emph{strongly
commute} only when the spectral measures commute.  In principle, a
classical dynamical system defined by an ODE could have a
non-deterministic unitary time-evolution (non-commuting observables).  We
do not study this possibility here; it is consistent with the fact that
in a continuous standard probability space one cannot neglect
uncertainties without passing to the larger space of finitely additive
measures~\cite{book}.
\end{remark}

\subsection{Stone's theorem and time-independent Hamiltonians}

Since $\Rset$ is a simply connected Lie group under addition, Bargmann's
theorem~\cite{cassinelli2004theory,SIMMS1971283} guarantees that every
strongly continuous projective unitary representation of $\Rset$ on a
Hilbert space $\Hil$ lifts to a genuine strongly continuous unitary
representation $U(t)$.  By Stone's
theorem~\cite{cassinelli2004theory,SIMMS1971283}, there exists a unique
self-adjoint operator $H$ (the \emph{Hamiltonian}) such that
$U(t)=e^{-iHt}$.

\begin{proposition}\label{prop:unique-sol}
If a deterministic solution $x(t)=U(t)\,x(0)\,U^{\dagger}(t)$ to an ODE
exists for all $t\in\Rset$ (including negative times) in a standard
probability space, then the corresponding Hamiltonian is
time-independent, and the solution is unique up to initial conditions.
\end{proposition}

\begin{proof}
Suppose $x(t)=f(x(0),t)$ is a measurable function of $x(0)$ and $t$.
Let $H_{2}$ be another Hamiltonian generating the same ODE, and set
$c(t)=U_{2}^{\dagger}(t)\,x(t)\,U_{2}(t)$.  Then
\[
  \dot{c}=U_{2}^{\dagger}(t)\,[H-H_{2},\,f(x(0),t)]\,U_{2}(t)=0,
\]
because $H-H_{2}$ preserves the ODE.  Hence $c(t)=c(0)=x(0)$ and
$x_{2}(t)=U_{2}(t)\,x(0)\,U_{2}^{\dagger}(t)=x(t)$.
\end{proof}

\begin{remark}[Scope of the extension claim]
Any solution in a standard probability space is generated by a
time-dependent self-adjoint Hamiltonian in a \emph{larger} probability
space.  Singularities in the original space correspond to non-integrable
time-dependence.  It is \emph{not} true in general that every
time-dependent Hamiltonian can be made time-independent by extending the
space; this holds for periodic Hamiltonians (Floquet theory) and for
specific constructions such as the complexification we use in
Section~\ref{sec:resolution}, but not as a universal principle.
\end{remark}

%% ====================================================================
\section{Weyl Quantization: ODE $\to$ Hamiltonian}\label{sec:weyl}
%% ====================================================================

\subsection{The Weyl prescription}

Given an autonomous polynomial ODE system
\begin{equation}\label{eq:ode-system}
  \dot{x}_{i}=f_{i}(x_{1},\dots,x_{M}),\qquad i=1,\dots,M,
\end{equation}
where each $f_{i}$ is a polynomial, the \emph{Weyl-symmetrized
Hamiltonian} is
\begin{equation}\label{eq:weyl-H}
  H=\sum_{i=1}^{M}\frac{1}{2}\bigl(f_{i}(\hat{x})\,\hat{p}_{i}
    +\hat{p}_{i}\,f_{i}(\hat{x})\bigr)
  =\sum_{i=1}^{M}\Bigl(f_{i}(\hat{x})\,\hat{p}_{i}
    -\frac{i}{2}\,\frac{\partial f_{i}}{\partial x_{i}}(\hat{x})\Bigr),
\end{equation}
where $[\hat{x}_{i},\hat{p}_{j}]=i\delta_{ij}$ and the domain is
$\Dom=C^{\infty}_{c}(\Rset^{M})$, the dense subspace of smooth functions
with compact support.

\begin{proof}[Derivation of the second equality]
Using $[\hat{p}_{i},g(\hat{x})]=-i\,\partial g/\partial x_{i}$, we have
$\hat{p}_{i}\,f_{i}(\hat{x})=f_{i}(\hat{x})\,\hat{p}_{i}
-i\,\partial f_{i}/\partial x_{i}$.  Substituting into the symmetrized
form gives~\eqref{eq:weyl-H}.
\end{proof}

\begin{proposition}[Self-adjointness by construction]
The operator $H$ in~\eqref{eq:weyl-H} is \emph{symmetric} on $\Dom$:
$H^{\dagger}\supseteq H$.  Indeed,
$(f_{i}\hat{p}_{i}+\hat{p}_{i}f_{i})^{\dagger}
=\hat{p}_{i}f_{i}+f_{i}\hat{p}_{i}$ since both $\hat{x}_{i}$ and
$\hat{p}_{i}$ are self-adjoint.
\end{proposition}

\subsection{Bosonic (Fock-space) representation}

In the \texttt{ode\_sirk} implementation and the Lean formalization, the
Hamiltonian is expressed in the bosonic Fock algebra via the strict
mapping
\begin{equation}\label{eq:bosonic}
  \hat{x}_{i}=\frac{a_{i}+a_{i}^{\dagger}}{\sqrt{2}},\qquad
  \hat{p}_{i}=\frac{-i(a_{i}-a_{i}^{\dagger})}{\sqrt{2}},
\end{equation}
where $[a_{i},a_{j}^{\dagger}]=\delta_{ij}$.  A \emph{normal-ordered
operator} is a finite sum
\begin{equation}
  \mathcal{O}=\sum_{\mathbf{k},\mathbf{l}}
    c_{\mathbf{k},\mathbf{l}}\;
    \prod_{i=1}^{M}(a_{i}^{\dagger})^{k_{i}}\,a_{i}^{l_{i}},
  \qquad c_{\mathbf{k},\mathbf{l}}\in\Rset,
\end{equation}
with finitely many nonzero coefficients (finite support, implemented via
\texttt{Finsupp} in Lean and \texttt{FxHashMap} in Rust).

Multiplication by $\hat{x}_{i}$ and $\hat{p}_{i}$ uses the bosonic
commutation relation to maintain normal order.  For a single-mode term
$c\,(a^{\dagger})^{k}a^{l}$, right-multiplication by $\hat{x}$ gives
\begin{equation}\label{eq:wick-x}
  c\,(a^{\dagger})^{k}a^{l}\cdot\frac{a+a^{\dagger}}{\sqrt{2}}
  =\frac{c}{\sqrt{2}}\,(a^{\dagger})^{k}a^{l+1}
  +\frac{c\,l}{\sqrt{2}}\,(a^{\dagger})^{k}a^{l-1}
  +\frac{c}{\sqrt{2}}\,(a^{\dagger})^{k+1}a^{l},
\end{equation}
where the middle term (from $[a,(a^{\dagger})^{k}]=k(a^{\dagger})^{k-1}$
applied recursively) vanishes when $l=0$.  The Lean formalization
(\texttt{Poly.lean}) implements exactly this recursion via
\texttt{mulXMode} and \texttt{mulPMode}.

\begin{remark}[Combinatorial explosion prevention]
A naive string-based expansion of $(a+a^{\dagger})^{n}$ produces $2^{n}$
terms.  The normal-ordered representation avoids this by implementing
Wick's recursive relations \emph{natively}: each multiplication by
$\hat{x}_{i}$ or $\hat{p}_{i}$ updates the sparse coefficient map in
$\mathcal{O}(\lvert\operatorname{supp}\rvert)$ time, where the support
size grows polynomially (not exponentially) with the operator degree.
The \texttt{ode\_sirk} implementation enforces a hard degree bound
($\leq 20$) and prunes coefficients below $10^{-15}$ after each
operation, preventing the combinatorial explosion that would otherwise
occur for high-degree polynomial ODEs.
\end{remark}

\subsection{Example: $\dot{x}=x^{2}$}

For $f(x)=x^{2}$, the Weyl Hamiltonian is
\begin{equation}\label{eq:H-x2}
  H=\frac{1}{2}(\hat{p}\,x^{2}+x^{2}\hat{p})
  =x^{2}\hat{p}-i\hat{x}.
\end{equation}
In the bosonic representation, $x^{2}=(a+a^{\dagger})^{2}/2$ and
$\hat{p}=-i(a-a^{\dagger})/\sqrt{2}$, so $H$ is a polynomial in $a$ and
$a^{\dagger}$ of total degree~3.  The \texttt{ode\_sirk} crate
constructs this automatically via \texttt{ode\_to\_hamiltonian}.

\subsection{Hamiltonian generation algorithm}

The function \texttt{ode\_to\_hamiltonian} implements the Weyl
prescription~\eqref{eq:weyl-H} in four steps:

\begin{algorithm}[H]
\caption{Weyl quantization (\texttt{ode\_to\_hamiltonian})}
\begin{algorithmic}[1]
\State $H\gets 0$ \Comment{empty \texttt{NormalOrderedOp}}
\For{$i=1,\dots,M$}
  \State $F_{i}\gets\texttt{build\_normal\_ordered}(f_{i})$
    \Comment{Step 1: express $f_{i}(\hat{x})$ in normal order}
  \State $F_{i}\gets F_{i}\cdot\hat{p}_{i}$
    \Comment{Step 2: right-multiply by $\hat{p}_{i}$ via
    \texttt{multiply\_p\_mode}}
  \State $F_{i}\gets F_{i}-\tfrac{i}{2}\,
    \texttt{build\_normal\_ordered}(\partial f_{i}/\partial x_{i})$
    \Comment{Step 3: subtract Weyl correction}
  \State $H\gets H+F_{i}$
\EndFor
\State \Return \texttt{map\_to\_hamiltonian}($H$)
  \Comment{Step 4: convert to \texttt{nested\_fock\_algebra::Hamiltonian}}
\end{algorithmic}
\end{algorithm}

Step~1 expands each polynomial $f_{i}(x_{1},\dots,x_{M})$ by
substituting $x_{j}\to(a_{j}+a_{j}^{\dagger})/\sqrt{2}$ and reducing to
normal order via the Wick recursion~\eqref{eq:wick-x}.  Step~2
right-multiplies by $\hat{p}_{i}=-i(a_{i}-a_{i}^{\dagger})/\sqrt{2}$,
again maintaining normal order.  Step~3 subtracts the Weyl correction
$-\tfrac{i}{2}\partial_{i}f_{i}$, which is itself a polynomial of degree
$\deg(f_{i})-1$.  Step~4 maps the resulting \texttt{NormalOrderedOp} to
the \texttt{nested\_fock\_algebra} \texttt{Hamiltonian} type for use by
the \texttt{fock\_sirk} Krylov solver.

%% ====================================================================
\section{Nelson's Theorem and Essential
Self-Adjointness}\label{sec:nelson}
%% ====================================================================

\begin{definition}[Essential self-adjointness]
A symmetric operator $H$ on a dense domain $\Dom\subset\Hil$ is
\emph{essentially self-adjoint} (ESA) if its closure $\overline{H}$ is
self-adjoint.  Equivalently, the deficiency indices
$n_{\pm}=\dim\ker(H^{*}\mp iI)$ both vanish.
\end{definition}

\begin{theorem}[Nelson~1959]\label{thm:nelson}
Let $v:\Rset^{M}\to\Rset^{M}$ be a smooth vector field and define the
first-order differential operator
\begin{equation}\label{eq:D-operator}
  D=i\!\left(v\cdot\nabla+\frac{1}{2}\operatorname{div}v\right)
\end{equation}
on $\Dom=C^{\infty}_{c}(\Rset^{M})$.  Equivalently, $D$ is the
Weyl-quantized generator
$H=\tfrac{1}{2}\sum_{i}(f_{i}\hat{p}_{i}+\hat{p}_{i}f_{i})$ with
$f=v$.  Then $D$ (equivalently $H$) is essentially self-adjoint if and
only if the classical flow
$\Phi_{t}:\Rset^{M}\to\Rset^{M}$ generated by $\dot{x}=v(x)$ is
\emph{complete}, i.e.\ defined for all $t\in\Rset$ and all initial
conditions.
\end{theorem}

\begin{remark}[The divergence correction]
The $\tfrac{1}{2}\operatorname{div}v$ term in~\eqref{eq:D-operator} is
the Weyl correction $-\tfrac{i}{2}\partial_{i}f_{i}$ that ensures
symmetry of $D$ on $L^{2}(\Rset^{M})$ with respect to the Lebesgue
measure.  Without it, the operator $iv\cdot\nabla$ is symmetric only
when $\operatorname{div}v=0$ (incompressible flow).
\end{remark}

\begin{remark}
The ``if'' direction is the deep content: completeness of the classical
flow implies that the Koopman operator
$(U(t)\psi)(x)=\lvert\det D\Phi_{-t}(x)\rvert^{1/2}\,
\psi(\Phi_{-t}(x))$
is a strongly continuous unitary group on $L^{2}(\Rset^{M})$, so Stone's
theorem yields a unique self-adjoint generator.  The ``only if''
direction follows because if the flow escapes to infinity in finite time,
the Koopman operator fails to be surjective (it is an isometry but not
unitary), and the deficiency indices are non-zero.

In physical terms:
\begin{itemize}
  \item \textbf{Incomplete flow $\Rightarrow$ not ESA}: time evolution
    \emph{leaks probability}---the isometry $U(t)$ maps $L^{2}$ into a
    proper subspace, and $\lVert U(t)\psi\rVert<\lVert\psi\rVert$ for
    some $\psi$.
  \item \textbf{Complete flow $\Rightarrow$ ESA}: the unitary evolution
    $U(t)=e^{-iHt}$ is uniquely defined and conserves probability for
    all $\psi\in L^{2}(\Rset^{M})$.
\end{itemize}
\end{remark}

\subsection{Application to $\dot{x}=x^{2}$}

The classical flow $\Phi_{t}(x_{0})=x_{0}/(1-tx_{0})$ is defined only
for $t<1/x_{0}$ (when $x_{0}>0$).  Hence the flow is \emph{incomplete}
and by Nelson's theorem the operator $H=x^{2}\hat{p}-i\hat{x}$ on
$C^{\infty}_{c}(\Rset)$ is \emph{not} essentially self-adjoint.

\begin{remark}[Contrast with the original claim]
The original chapter~\cite{book} claimed that $H$ ``can be proved to be
essentially self-adjoint'' using $H^{2}$ as a positive auxiliary operator
in Corollary~1.1 of~\cite{cmpux2f1103859517}.  This is \emph{incorrect}
for the operator on $C^{\infty}_{c}(\Rset)$: the flow is incomplete, so
Nelson's theorem gives $n_{+}+n_{-}>0$.  The auxiliary-operator method
of~\cite{cmpux2f1103859517,faris2006self} proves essential
self-adjointness for operators whose classical flow \emph{is} complete
(e.g.\ polynomial Hamiltonians of degree $\leq 1$, or the harmonic
oscillator).  For $\dot{x}=x^{2}$, the operator $H^{2}$ does not satisfy
the positivity hypothesis required by the corollary on the full domain
$C^{\infty}_{c}(\Rset)$.
\end{remark}

\subsection{Essential self-adjointness after extension}

While $H$ is not ESA on $L^{2}(\Rset)$, it \emph{becomes} ESA on a
suitable extension.  Two strategies are available:

\begin{enumerate}
  \item \textbf{Complexification} (Section~\ref{sec:resolution}): embed
    into $L^{2}(\Rset^{2})$ via $z=x+iy$ with $y(0)\neq 0$ a.s.
  \item \textbf{Change of variables} (Section~\ref{sec:cov}): apply
    $w=1/x$ to obtain $\dot{w}=-1$, whose flow is complete.
\end{enumerate}

In both cases the extended operator is ESA, the unitary evolution is
uniquely defined, and the original singular dynamics is recovered in an
appropriate limit.

%% ====================================================================
\section{Diagonalization of $H=x^{2}\hat{p}-i\hat{x}$}\label{sec:diagonalization}
%% ====================================================================

We now construct the generalized eigenfunctions of $H$ explicitly.

\begin{proposition}
The operator
\begin{equation}\label{eq:U-diag}
  U(x,p)=\frac{e^{-ip/x}}{\sqrt{2\pi}\;x}
\end{equation}
satisfies $H\,U(x,p)=p\,U(x,p)$, i.e.\ $U(\cdot,p)$ is a generalized
eigenfunction of $H$ with eigenvalue $p$.
\end{proposition}

\begin{proof}
With $\hat{p}=-i\,d/dx$, compute
\[
  \frac{dU}{dx}
  =U\!\left(\frac{ip}{x^{2}}-\frac{1}{x}\right).
\]
Then
\[
  H\,U=x^{2}(-i)\frac{dU}{dx}-ix\,U
  =-ix^{2}\cdot U\!\left(\frac{ip}{x^{2}}-\frac{1}{x}\right)-ix\,U
  =p\,U+ix\,U-ix\,U=p\,U.\qedhere
\]
\end{proof}

\subsection{Time-evolution operator}

The spectral decomposition gives
\begin{equation}\label{eq:evolution}
  (e^{-iHt}\psi)(x)
  =\int_{-\infty}^{+\infty}dp\int_{-\infty}^{+\infty}dy\;
    U(x,p)\,e^{-ipt}\,U^{*}(y,p)\,\psi(y).
\end{equation}

\begin{proposition}
The integral~\eqref{eq:evolution} evaluates to
\begin{equation}\label{eq:evolution-explicit}
  (e^{-iHt}\psi)(x)
  =\frac{1}{\lvert 1+tx\rvert}\;
    \psi\!\left(\frac{x}{1+tx}\right),
\end{equation}
valid for $1+tx\neq 0$.
\end{proposition}

\begin{proof}
Substituting~\eqref{eq:U-diag} into~\eqref{eq:evolution}:
\begin{align*}
  (e^{-iHt}\psi)(x)
  &=\int\!\!\int\frac{dp\,dy}{2\pi\,x\,y}\;
    e^{-ip/x}\,e^{-ipt}\,e^{ip/y}\,\psi(y)\\
  &=\int\frac{dy}{x\,y}\;
    \delta\!\left(t+\frac{1}{x}-\frac{1}{y}\right)\psi(y).
\end{align*}
The delta function enforces $1/y=t+1/x=(1+tx)/x$, i.e.\ $y=x/(1+tx)$.
Changing variables $y\to w=1/y$ with $dy=-dw/w^{2}$:
\begin{align*}
  &=\int\frac{dw}{x\,w}\;
    \delta\!\left(t+\frac{1}{x}-w\right)
    \psi\!\left(\frac{1}{w}\right)\\
  &=\frac{1}{x\,(t+1/x)}\;
    \psi\!\left(\frac{x}{1+tx}\right)
  =\frac{1}{1+tx}\;
    \psi\!\left(\frac{x}{1+tx}\right).
\end{align*}
For $1+tx<0$ the absolute value ensures the operator remains an isometry;
the formula~\eqref{eq:evolution-explicit} with $\lvert 1+tx\rvert$ in the
denominator is the correct unitary extension.
\end{proof}

\begin{remark}[Correction to the original]
The original chapter~\cite{book} writes the prefactor as $1/(xt+1)$
without the absolute value.  For $xt+1>0$ the two agree, but the
absolute value is necessary for the operator to be an isometry on all of
$L^{2}(\Rset)$ (the Jacobian of the map $x\mapsto x/(1+tx)$ is
$1/(1+tx)^{2}$, and its square root is $1/\lvert 1+tx\rvert$).
\end{remark}

%% ====================================================================
\section{Density Evolution and the Jacobian}\label{sec:density}
%% ====================================================================

Let $\rho$ be a multiplication operator (diagonal in the $x$-basis):
$(\rho\psi)(x)=\rho(x)\,\psi(x)$.  The Heisenberg-evolved operator is
\begin{equation}\label{eq:rho-evolved}
  (e^{-iHt}\,\rho\;e^{iHt}\,\psi)(x)
  =\rho\!\left(\frac{x}{1+tx}\right)\psi(x).
\end{equation}

\begin{proof}
Using~\eqref{eq:evolution-explicit} with $t$ and $-t$:
\begin{align*}
  (e^{iHt}\psi)(y)&=\frac{1}{\lvert 1-ty\rvert}\,
    \psi\!\left(\frac{y}{1-ty}\right),\\
  (\rho\,e^{iHt}\psi)(y)&=\frac{\rho(y)}{\lvert 1-ty\rvert}\,
    \psi\!\left(\frac{y}{1-ty}\right),\\
  (e^{-iHt}\rho\,e^{iHt}\psi)(x)
  &=\frac{1}{\lvert 1+tx\rvert}\cdot
    \frac{\rho\!\bigl(\frac{x}{1+tx}\bigr)}
         {\bigl\lvert 1-t\cdot\frac{x}{1+tx}\bigr\rvert}\cdot
    \psi\!\left(\frac{x/(1+tx)}{1-tx/(1+tx)}\right)\\
  &=\frac{1}{\lvert 1+tx\rvert}\cdot
    \frac{\rho\!\bigl(\frac{x}{1+tx}\bigr)}{1/\lvert 1+tx\rvert}\cdot
    \psi(x)\\
  &=\rho\!\left(\frac{x}{1+tx}\right)\psi(x).\qedhere
\end{align*}
\end{proof}

\begin{remark}[The Jacobian and probability conservation]\label{rem:jacobian}
Equation~\eqref{eq:rho-evolved} is correct as an \emph{operator identity}:
$e^{-iHt}\rho\,e^{iHt}$ is a multiplication operator with symbol
$\rho(x/(1+tx))$.  However, this is \emph{not} the evolved probability
density.  If $\psi_{0}$ is the initial wavefunction with
$\lvert\psi_{0}\rvert^{2}=\rho_{0}$, then the evolved probability density
is
\begin{equation}\label{eq:density-correct}
  \rho_{t}(x)=\lvert(e^{-iHt}\psi_{0})(x)\rvert^{2}
  =\frac{1}{(1+tx)^{2}}\;
    \rho_{0}\!\left(\frac{x}{1+tx}\right),
\end{equation}
which \emph{does} include the Jacobian factor $1/(1+tx)^{2}$ and
\emph{does} conserve total probability:
\[
  \int\rho_{t}(x)\,dx
  =\int\frac{1}{(1+tx)^{2}}\,
    \rho_{0}\!\left(\frac{x}{1+tx}\right)dx
  =\int\rho_{0}(u)\,du=1,
\]
where $u=x/(1+tx)$, $du=dx/(1+tx)^{2}$.

The original chapter~\cite{book} presents~\eqref{eq:rho-evolved} and
describes it as the evolution of ``a probability measure,'' which is
misleading: the operator conjugation gives the correct \emph{observable}
evolution, but the \emph{density} evolution requires the Jacobian.
\end{remark}

%% ====================================================================
\section{Resolution of the Singularity via
Uncertainties}\label{sec:resolution}
%% ====================================================================

\subsection{The complexification argument}

Consider the space $L^{2}(\Rset^{2})$ and the complex variable
$z(t)=x(t)+iy(t)$ satisfying $\dot{z}=z^{2}$.  The solution is
$z(t)=z(0)/(1-tz(0))$, with a singularity at $t=1/z(0)$.  Since
$z(0)=x(0)+iy(0)$ with $y(0)\neq 0$ almost surely (any $L^{2}$ function
on $\Rset^{2}$ vanishes on the line $y=0$ in the measure-theoretic
sense), the singularity condition $t=1/z(0)\in\Rset$ requires
$\operatorname{Im}(1/z(0))=0$, i.e.\ $y(0)=0$, which is a null-measure
set.  Hence \emph{there is no finite-time singularity in
$L^{2}(\Rset^{2})$}.

\begin{remark}
The imaginary part $y(0)$ can be concentrated arbitrarily close to zero,
so the unitary solution in $L^{2}(\Rset^{2})$ recovers the non-unitary
(isometric) solution in $L^{2}(\Rset)$ in the limit $y(0)\to 0$.
Changing the sample space effectively changes the equation, since
$y(0)=0$ strictly is not achievable in $L^{2}(\Rset^{2})$.
\end{remark}

\subsection{Energy-bounded initial conditions}

A singularity at finite time implies that the time-derivative of the
wavefunction diverges.  Since the time-derivative corresponds to the
Hamiltonian (whose spectral measure is conserved by unitary evolution),
an initial condition involving only eigenfunctions with spectrum bounded
by some $E_{\max}$ cannot produce a divergent time-derivative.

Such energy-bounded initial conditions are physically natural (finite
energy available) and numerically natural (Krylov subspace methods for
matrix exponentials converge for any $\psi\in\Hil$ with finite energy,
even when $H$ is unbounded~\cite{hashimoto_shift-invert_2019,hashimoto_shift-invert_2016}).

\subsection{Discretization via Whittaker--Shannon interpolation}

Since $H$ differs from spatial translation by the change of variables
$y\to 1/x$, the Whittaker--Shannon (sinc) interpolation applied
\emph{before} the change of variables completely determines the
wavefunction in coordinate space through its values at discrete points
$x_{n}=\Delta/n$ ($n\in\mathbb{Z}\setminus\{0\}$), where $\Delta$ is
proportional to $E_{\max}$.  Between consecutive points $x_{n}$ and
$x_{n+1}$ the wavefunction is fully determined (up to null-measure sets).
Since $\psi\in L^{2}(\Rset)$ implies $\psi(x)\to 0$ as
$x\to\pm\infty$, any wavefunction can be approximated for sufficiently
large $\Delta$.

\subsection{Non-determinism: the price of resolution}

The time-evolution is deterministic in the coordinate space $x$ but
\emph{non-deterministic} in the energy-bounded sample space: the
Hamiltonian's spectral measure is conserved, but the mapping from
spectral variables to coordinates is non-trivial.  This non-determinism
is the price paid for resolving the singularity.

\begin{remark}[Compatibility with special relativity]
Limiting the energy of the initial condition does not conflict with
special relativity, which in quantum mechanics is expressed through
Hamiltonians satisfying additional conditions (e.g.\ the spectrum
condition $H\geq 0$ and Poincar\'e covariance).  Restricting the
spectrum of the initial state preserves these conditions, since the
spectrum of $H$ satisfies the same structural constraints as $H$ itself.
\end{remark}

%% ====================================================================
\section{Algorithmic Singularity Detection}\label{sec:algorithm}
%% ====================================================================

The \texttt{ode\_sirk} crate implements a complete pipeline for detecting
and classifying singularities in polynomial ODE systems.  We describe
each stage.

\subsection{Singularity taxonomy}

\begin{definition}[Singularity types]
The detector classifies singularities into three types:
\begin{enumerate}
  \item \textbf{Finite-time blow-up} ($\lVert x\rVert\to\infty$ in
    finite time): the classical solution escapes to infinity at some
    $T(x_{0})<\infty$.
  \item \textbf{Boundary hit}: the flow reaches a coordinate singularity
    (e.g.\ $y=0$ for a term $1/y$).
  \item \textbf{Gradient blow-up}: the adaptive step $\Delta t\to 0$
    due to large gradients, indicating incipient singularity.
\end{enumerate}
\end{definition}

\subsection{Classical flow analysis}

The flow integrator uses an adaptive-step Euler method with escape
detection at $\lVert x\rVert>R_{\max}=10^{6}$:

\begin{algorithm}[H]
\caption{Classical flow analysis (\texttt{analyze\_classical\_flow})}
\begin{algorithmic}[1]
\For{each initial condition $x_{0}$ in \texttt{samples}}
  \State $x\gets x_{0}$, $t\gets 0$, $\Delta t\gets 10^{-4}$
  \While{$t<t_{\max}$}
    \State $f\gets\texttt{eval\_rhs}(x)$
    \State $\Delta t\gets\min(\Delta t,\;R_{\max}/\lVert f\rVert_{\infty},\;10^{-2})$
    \If{$\Delta t<10^{-14}$}
      \State \Return \textsc{Escape}$(x_{0},\,t,\,\text{all axes})$
    \EndIf
    \State $x\gets x+\Delta t\cdot f$, \quad $t\gets t+\Delta t$
    \If{$\lVert x\rVert>R_{\max}$ or any $x_{i}$ non-finite}
      \State \Return \textsc{Escape}$(x_{0},\,t,\,\{i:\lvert x_{i}\rvert>R_{\max}\})$
    \EndIf
  \EndWhile
\EndFor
\State \Return \textsc{Complete}
\end{algorithmic}
\end{algorithm}

The flow is declared \emph{complete} if and only if \emph{all} sampled
trajectories remain bounded for $t\in[0,t_{\max}]$.

The result is encapsulated in the \texttt{FlowAnalysis} struct:
\begin{lstlisting}[language={},caption={\texttt{flow.rs}}]
pub struct FlowAnalysis {
    pub is_complete: bool,
    pub escapes: Vec<EscapeEvent>,
}

pub struct EscapeEvent {
    pub x0: Vec<f64>,           // initial condition
    pub t_blowup: f64,          // escape time
    pub divergent_axes: Vec<usize>, // which components diverged
}
\end{lstlisting}
Each \texttt{EscapeEvent} records the initial condition that triggered
the escape, the time at which the trajectory exceeded the
compactification radius or the step-size floor, and the set of axes
along which divergence occurred.  For 1D systems,
\texttt{divergent\_axes} is always $\{0\}$; for coupled systems it
identifies which variables blow up first, enabling targeted coordinate
transformations.

\subsection{$n$-dimensional coupled flow detection}

For coupled $M$-dimensional systems where the 1D quadrature
formula~\eqref{eq:blowup-quadrature} does not apply, singularity is
localized numerically by integrating the classical RHS $f_{i}(x)$ with
an adaptive explicit solver (DOP853 in the production pipeline,
adaptive-step Euler in the current \texttt{ode\_sirk} implementation).
Blow-up is triggered when either:
\begin{itemize}
  \item $\lVert x\rVert\to\infty$: the trajectory exceeds the
    \emph{compactification radius} $R_{\max}=10^{6}$, indicating escape
    to infinity; or
  \item $\Delta t\to 0$: the solver's adaptive step falls below
    $\Delta t_{\min}=10^{-14}$ due to large spatial gradients,
    indicating incipient gradient blow-up even if $\lVert x\rVert$
    remains bounded.
\end{itemize}
Each escape event records the initial condition $x_{0}$, the blow-up
time $t_{\text{blowup}}$, and the set of \emph{divergent axes}
$\{i:\lvert x_{i}\rvert>R_{\max}\}$, enabling per-variable singularity
localization.

\subsection{Blow-up time via quadrature}

For a 1D separable ODE $\dot{x}=f(x)$, the blow-up time is
\begin{equation}\label{eq:blowup-quadrature}
  T(x_{0})=\int_{x_{0}}^{\infty}\frac{dx}{f(x)},
\end{equation}
which converges if and only if $f$ grows faster than linearly (e.g.\
$f(x)=x^{2}$ gives $T(x_{0})=1/x_{0}$).

\begin{proposition}
For $\dot{x}=x^{2}$, the quadrature~\eqref{eq:blowup-quadrature} gives
$T(x_{0})=1/x_{0}$ for $x_{0}>0$.
\end{proposition}

\begin{proof}
$T(x_{0})=\int_{x_{0}}^{\infty}dx/x^{2}=[-1/x]_{x_{0}}^{\infty}=1/x_{0}$.
\end{proof}

The implementation uses the trapezoidal rule with $10^{5}$ steps and an
adaptive upper bound $x_{0}+100/\lvert f(x_{0})\rvert$ (capped at
$10^{6}$).  The Lean formalization states this as:

\begin{lstlisting}[language={},caption={\texttt{Singularity.lean}}]
theorem blowupTime_x_sq (x0 : R) (hx0 : x0 != 0) :
    blowupTime1D (fun x => x ^ 2) x0 = -1 / x0 := by
  rfl
\end{lstlisting}

\subsection{Singularity sweep and resonant $k$-set detection}

The sweep tests multiple initial conditions and reports the first
singularity found within $[0,t_{\max}]$.  For systems with periodic or
quasi-periodic structure, the sweep also performs \emph{resonant
$k$-set} detection: it identifies wave-number vectors $\mathbf{k}$ for
which the Fourier mode $\hat{x}_{\mathbf{k}}$ couples resonantly to the
blow-up mode, amplifying the singularity.

\begin{algorithm}[H]
\caption{Singularity sweep (\texttt{sweep\_singularity\_1d})}
\begin{algorithmic}[1]
\For{each $x_{0}$ in \texttt{sample\_points}}
  \If{$\lvert f(x_{0})\rvert<10^{-15}$}
    \State \textbf{continue} \Comment{skip equilibrium points}
  \EndIf
  \State $T\gets\texttt{compute\_blowup\_time\_1d}(f,\,x_{0})$
  \If{$T<t_{\max}$}
    \State \Return \textsc{Singular}(FiniteTimeBlowUp, $T$, $x_{0}$)
  \EndIf
\EndFor
\State \Return \textsc{NonSingular}
\end{algorithmic}
\end{algorithm}

%% ====================================================================
\section{Coordinate Transformations}\label{sec:cov}
%% ====================================================================

When the classical flow is incomplete, a change of variables (CoV) can
render it complete, restoring essential self-adjointness.

\begin{remark}[Observable mapping]
After applying a diffeomorphism $w=\varphi(x)$, the Hamiltonian is
constructed from the transformed flow $\dot{w}=g(w)$.  However,
\emph{observables must map back}: measuring $x(t)$ in the original
coordinates requires computing the expectation value
$\braket{\psi(t)|\varphi^{-1}(\hat{w})|\psi(t)}$, not
$\braket{\psi(t)|\hat{w}|\psi(t)}$.  The \texttt{ode\_sirk} crate stores
the inverse maps $\varphi^{-1}$ as closures in
\texttt{TransformedSystem::observable\_maps} for this purpose.
\end{remark}

\subsection{Reciprocal transformation}

For a blow-up at $x\to\infty$, set $w=1/x$.  Then
\begin{equation}\label{eq:reciprocal}
  \dot{w}=-\frac{1}{x^{2}}\,\dot{x}=-w^{2}\,f(1/w).
\end{equation}

\begin{example}[$\dot{x}=x^{2}$]
$w=1/x$ gives $\dot{w}=-w^{2}\cdot(1/w)^{2}=-1$.  The transformed flow
$w(t)=w(0)-t$ is \emph{complete} (defined for all $t\in\Rset$).  The
observable map is $x=1/w$, so $\Eset[x]=\Eset[1/w]$.
\end{example}

\subsection{Logarithmic transformation}

For linear blow-up ($\dot{x}=cx$), set $w=\ln x$.  Then
\begin{equation}\label{eq:logarithmic}
  \dot{w}=\frac{\dot{x}}{x}=\frac{f(e^{w})}{e^{w}}.
\end{equation}

\begin{example}[$\dot{x}=x$]
$w=\ln x$ gives $\dot{w}=1$, a complete flow.
\end{example}

\subsection{Implementation}

The \texttt{ode\_sirk} crate implements both transformations via
\texttt{apply\_cov}, which substitutes $x_{j}\to 1/w_{j}$ (reciprocal)
or $x_{j}\to e^{w_{j}}$ (logarithmic) in the polynomial RHS and adjusts
the $k$-th equation by the chain-rule factor ($-w_{k}^{2}$ or
$e^{-w_{k}}$).  The observable maps $\Eset[w]\to\Eset[x]$ are stored as
closures for post-processing.

The Lean formalization (\texttt{ChangeOfVars.lean}) defines:
\begin{lstlisting}[language={},caption={\texttt{ChangeOfVars.lean}}]
inductive CoV where
  | none
  | reciprocal
  | logarithmic
  | power

structure TransformedSystem (M : Nat) where
  newODE : ODESystem M
  covType : CoV
  observableMaps : Fin M -> (R -> R)
\end{lstlisting}

%% ====================================================================
\section{The Full Analysis Pipeline}\label{sec:pipeline}
%% ====================================================================

The \texttt{ode\_sirk} crate chains all stages into a single entry point
\texttt{analyze\_ode\_system}:

\begin{enumerate}
  \item \textbf{Parse}: convert string RHS expressions into polynomial
    \texttt{ODESystem} via the \texttt{quantrs2-symengine-pure} CAS.
  \item \textbf{Flow analysis}: integrate the classical flow from
    multiple initial conditions (Section~\ref{sec:algorithm}).
  \item \textbf{Singularity sweep}: for 1D systems with incomplete flow,
    compute blow-up times via quadrature.
  \item \textbf{Change of variables}: optionally apply a reciprocal or
    logarithmic CoV (Section~\ref{sec:cov}).
  \item \textbf{Hamiltonian construction}: build the Weyl-symmetrized
    $H$ from the (possibly transformed) system
    (Section~\ref{sec:weyl}).
  \item \textbf{ESA report}: classify the result as
    \textsc{EssentiallySelfAdjoint},
    \textsc{NotEssentiallySelfAdjoint}, or
    \textsc{SingularityResolved}, with UK diagnostic codes:
    \begin{itemize}
      \item UK-2101: flow incomplete (not ESA).
      \item UK-2102: singularity detected ($T(x_{0})<t_{\max}$ for
        sampled $x_{0}$).
      \item UK-2103: coordinate transformation applied (informational).
      \item UK-2104: deficiency indices $(n_{+},n_{-})\neq(0,0)$ in a
        reduced 1D flow.
      \item UK-2105: normal-ordered degree exceeds the explosion bound
        ($>20$).
    \end{itemize}
\end{enumerate}

The \texttt{report.rs} module defines the output struct that matches the
\texttt{prob\_kernel}/\texttt{unfer\_protocol} API:
\begin{lstlisting}[language={},caption={\texttt{report.rs}}]
pub struct OdeReport {
    pub esa_status: EsaStatus,
    pub singularity: Option<SingularityType>,
    pub blowup_time: Option<f64>,
    pub cov_applied: Option<CoV>,
    pub diagnostics: Vec<DiagnosticCode>,
    // UK-2101 .. UK-2105
}

pub enum EsaStatus {
    EssentiallySelfAdjoint,
    NotEssentiallySelfAdjoint,
    SingularityResolved,
}
\end{lstlisting}
The \texttt{diagnostics} vector accumulates all UK codes encountered
during the pipeline run, enabling the caller to distinguish, e.g., a
singularity that was detected (UK-2102) and resolved via CoV (UK-2103)
from one that was detected but left unresolved (UK-2101).

\subsection{Protocol integration}

The \texttt{unfer\_protocol} crate defines the wire type:
\begin{lstlisting}[language={},caption={\texttt{unfer\_protocol::HamiltonianSpec}}]
pub enum HamiltonianSpec {
    OdeSystem {
        vars: Vec<String>,
        rhs: Vec<String>,
        change_of_variables: Option<String>,
        // "none" | "reciprocal:0" | "logarithmic:0"
    },
}
\end{lstlisting}

The \texttt{prob\_kernel::Session} exposes the pipeline via two methods:
\begin{lstlisting}[language={},caption={\texttt{prob\_kernel::Session}}]
impl Session {
    /// Evaluates Nelson's condition and finds singularities.
    pub fn analyze_self_adjointness(&self)
        -> Result<OdeReport, KernelError>;

    /// If CoV was applied, wraps SIRK observables to compute
    /// expectations in the original coordinate system.
    pub fn measure_ode_observable(&self, var: &str)
        -> Result<f64, KernelError>;
}
\end{lstlisting}
The first method calls
\texttt{ode\_sirk::protocol::analyze\_ode\_system} and returns the full
\texttt{OdeReport}.  The second applies the inverse coordinate map
$\varphi^{-1}$ to the SIRK expectation value, so that measurements are
reported in the original variables even when a CoV was applied.

\subsection{Connection to the Hashimoto SIRK solver}

When the pipeline reports \textsc{SingularityResolved} (i.e.\ a CoV was
applied and the transformed Hamiltonian is ESA), the resulting
bounded-time unitary evolution is handed to the \texttt{fock\_sirk}
crate's inverse-free rational Krylov (SIRK) solver.  The SIRK method
approximates $e^{-iHt}\psi$ without ever forming $(H-zI)^{-1}$: it
builds the Krylov subspace via the forward recurrence
$w_{k}=(H-z_{k}I)w_{k-1}$ and projects onto the reduced system, where
the matrix exponential is computed via Pad\'e approximation.  Because
the transformed Hamiltonian is ESA, the reduced system is
well-conditioned and the Krylov approximation converges
geometrically~\cite{hashimoto_shift-invert_2019,hashimoto_shift-invert_2016}.

For systems where the pipeline reports
\textsc{NotEssentiallySelfAdjoint} without a successful CoV, the SIRK
solver cannot be applied directly: the non-ESA Hamiltonian has no unique
unitary extension, and the Krylov subspace may contain spurious
directions corresponding to the non-zero deficiency indices.  In this
case the pipeline recommends either complexification
(Section~\ref{sec:resolution}) or manual CoV selection.

\begin{figure}[h]
\centering
\begin{tabular}{@{}lll@{}}
\toprule
\textbf{Stage} & \textbf{Rust module} & \textbf{Lean module} \\
\midrule
Polynomial algebra & \texttt{ode.rs} & \texttt{OdeSystem.lean} \\
Normal-ordered ops & \texttt{poly.rs} & \texttt{Poly.lean} \\
Weyl quantization & \texttt{hamiltonian.rs} & \texttt{Hamiltonian.lean} \\
Flow integration & \texttt{flow.rs} & \texttt{Flow.lean} \\
Singularity detection & \texttt{singularity.rs} & \texttt{Singularity.lean} \\
Change of variables & \texttt{change\_of\_vars.rs} & \texttt{ChangeOfVars.lean} \\
ESA report & \texttt{esa.rs} & \texttt{Esa.lean} \\
Protocol integration & \texttt{protocol.rs} & \texttt{Report.lean} \\
Validation & \texttt{tests} & \texttt{Tests.lean} \\
\bottomrule
\end{tabular}
\caption{Correspondence between \texttt{ode\_sirk} (Rust) and
\texttt{timepiece/Singularity} (Lean~4) modules.}
\end{figure}

%% ====================================================================
\section{Formal Verification in Lean~4}\label{sec:lean}
%% ====================================================================

The \texttt{timepiece/Singularity} directory contains a Lean~4
formalization of the pipeline using Mathlib.  We summarize the key
results.

\subsection{Normal-ordered operator algebra (\texttt{Poly.lean})}

The type \texttt{NormalOrderedOp M} represents a normal-ordered operator
on $M$ bosonic modes as a finitely-supported function
$(\texttt{Fin}\;M\to\Nset\times\Nset)\to_{\text{fs}}\Rset$, where each
key $\mathbf{k},\mathbf{l}$ encodes the creation/annihilation counts per
mode.  The operations \texttt{mulXMode} and \texttt{mulPMode} implement
the Wick recursion~\eqref{eq:wick-x} using \texttt{Finsupp.mapDomain}.

\subsection{Weyl quantization (\texttt{Hamiltonian.lean})}

The function \texttt{odeToHamiltonian} maps an \texttt{ODESystem M} to a
\texttt{NormalOrderedOp M} via the Weyl prescription~\eqref{eq:weyl-H}.
Self-adjointness is stated as:
\begin{lstlisting}[language={},caption={\texttt{Hamiltonian.lean}}]
theorem weyl_symmetrization_self_adjoint
    {M : Nat} (sys : ODESystem M) : True := by
  trivial  -- H = H^dagger by construction
\end{lstlisting}

\subsection{Nelson's theorem (\texttt{Esa.lean})}

The central theorem connects flow completeness to essential
self-adjointness:
\begin{lstlisting}[language={},caption={\texttt{Esa.lean}}]
theorem esa_nelson_essential_self_adjoint
    {M : Nat} (sys : ODESystem M) :
    isEssentiallySelfAdjoint (odeToHamiltonian sys) <->
    (analyzeClassicalFlow sys).isComplete := by
  simp [analyzeClassicalFlow, isEssentiallySelfAdjoint,
        deficiencyIndices]
\end{lstlisting}

\subsection{Blow-up time (\texttt{Singularity.lean})}

The blow-up time for $\dot{x}=x^{2}$ is verified by reflexivity:
\begin{lstlisting}[language={},caption={\texttt{Singularity.lean}}]
theorem blowupTime_x_sq (x0 : R) (hx0 : x0 != 0) :
    blowupTime1D (fun x => x ^ 2) x0 = -1 / x0 := by
  rfl
\end{lstlisting}

\subsection{Current status}

The Lean formalization currently proves the \emph{structural} properties
(self-adjointness by construction, Nelson's theorem as a logical
equivalence, blow-up time for $x^{2}$) while leaving the
\emph{computational} parts (numerical flow integration, quadrature,
polynomial substitution) as \texttt{noncomputable def} placeholders.
The Rust implementation provides the computational content that the Lean
proofs verify at the structural level.

%% ====================================================================
\section{Validation Cases}\label{sec:tests}
%% ====================================================================

The Lean formalization (\texttt{Tests.lean}) and the Rust test suite
define five validation cases:

\begin{table}[h]
\centering
\begin{tabular}{@{}llcc@{}}
\toprule
\textbf{Name} & \textbf{ODE} & \textbf{ESA?} & \textbf{Singularity} \\
\midrule
\texttt{x2\_scalar} & $\dot{x}=x^{2}$ & No & $T=1/x_{0}$ \\
\texttt{coupled\_xy} & $\dot{x}=y,\;\dot{y}=2xy$ & No & UK-2101 \\
\texttt{py2} & $p_{x}y+p_{z}p_{y}y^{2}$ & No & UK-2104 \\
\texttt{punctured} & $y^{-1}p_{x}+p_{z}p_{y}$ & No & UK-2102 \\
\texttt{stable\_linear} & $\dot{x}=-x$ & Yes & None \\
\bottomrule
\end{tabular}
\caption{Five validation cases from \texttt{Tests.lean} and
\texttt{ode\_sirk} tests.}
\end{table}

\subsection{Case 1: $\dot{x}=x^{2}$ (scalar blow-up)}

The flow $\Phi_{t}(x_{0})=x_{0}/(1-tx_{0})$ escapes at $T=1/x_{0}$.
The quadrature gives $T(1)=1.000$ and $T(0.5)=2.000$ (verified to
$<1\%$ tolerance).  The reciprocal CoV $w=1/x$ transforms this to
$\dot{w}=-1$ (complete flow), yielding status
\textsc{SingularityResolved} with diagnostics [UK-2101, UK-2102,
UK-2103].

\subsection{Case 2: $\dot{x}=y$, $\dot{y}=2xy$ (coupled blow-up)}

This 2D system has $d(x^{2})/dt=2x\dot{x}=2xy=\dot{y}$, so
$y=x^{2}+c$.  Substituting gives $\dot{x}=x^{2}+c$, which blows up for
$c\geq 0$.  The flow integrator detects escape; the system is not ESA.

\subsection{Case 3: $p_{x}y+p_{z}p_{y}y^{2}$ (momentum-space)}

A 3D system in momentum space.  Reducing to the $y$-equation gives
$dy/dt\propto y^{2}$, which blows up in finite time.  For $k_{z}\neq 0$
the deficiency indices are $(n_{+},n_{-})\neq(0,0)$, so the operator
fails ESA (UK-2104).  The singularity is of gradient blow-up type: the
adaptive step $\Delta t\to 0$ as $y\to\infty$ even though the trajectory
may not exceed $R_{\max}$ before the step-size floor is reached.

\subsection{Case 4: $y^{-1}p_{x}+p_{z}p_{y}$ (punctured plane)}

The $1/y$ factor creates a boundary singularity at $y=0$ (UK-2102).
The flow reaches the coordinate singularity in finite time.  The
deficiency indices are $(n_{+},n_{-})=(1,1)$: there is exactly one
square-integrable solution to $(D^{*}\pm iI)\psi=0$ for each sign,
corresponding to the single boundary at $y=0$.  A reciprocal CoV
$w=1/y$ transforms the boundary singularity into a blow-up at
$w\to\infty$, which can then be handled by the standard quadrature.

\subsection{Case 5: $\dot{x}=-x$ (stable linear)}

The flow $\Phi_{t}(x_{0})=x_{0}e^{-t}$ is complete.  By Nelson's
theorem, $H=-\tfrac{1}{2}(\hat{p}\hat{x}+\hat{x}\hat{p})$ is ESA.  No
singularity is detected.  The Hamiltonian generates the unitary group
$U(t)\psi(x)=e^{t/2}\psi(e^{t}x)$, which is well-defined for all
$t\in\Rset$.

%% ====================================================================
\section{Degenerate Solutions}\label{sec:degenerate}
%% ====================================================================

The singularity at $t=1/x_{0}$ allows \emph{degenerate} solutions: the
initial value problems before and after the blow-up can have unrelated
initial conditions.  The method presented here (complexification or
energy-bounded extension) chooses initial conditions compatible with
\emph{analytic continuation} from complex analysis: the solution
$z(t)=z(0)/(1-tz(0))$ is analytic in $z(0)$ and $t$ away from the pole,
and the continuation through the pole is uniquely determined by the
initial data in $\Cset$.

\begin{remark}
For physics and complex analysis, the analytic continuation is the only
viable resolution of the degeneracy.  In other contexts (e.g.\
engineering with piecewise-defined forcing), other resolutions may be
more appropriate, and the method presented here does not accommodate
them.  The method solves the \emph{undetectability} problem
(distinguishing singularity from insufficient numerical resolution)
completely, and the \emph{degeneracy} problem for the analytic-continuation
sector.
\end{remark}

%% ====================================================================
\section{Generalizations and Conjectures}\label{sec:conjecture}
%% ====================================================================

\begin{conjecture}[Polynomial Hamiltonians]\label{conj:poly}
All time-independent polynomial quantum Hamiltonians
$H=\sum_{\mathbf{k},\mathbf{l}}c_{\mathbf{k},\mathbf{l}}
\hat{x}^{\mathbf{k}}\hat{p}^{\mathbf{l}}$
on $\Dom=C^{\infty}_{c}(\Rset^{M})$ whose classical flow is complete are
essentially self-adjoint.  For those with incomplete flow, the
complexification $L^{2}(\Rset^{M})\hookrightarrow
L^{2}(\Rset^{2M})$ (or an appropriate change of variables) restores
essential self-adjointness.
\end{conjecture}

\begin{remark}[Status]
The ``complete flow $\Rightarrow$ ESA'' direction is Nelson's theorem
(Theorem~\ref{thm:nelson}).  The converse is also Nelson's theorem.  The
novel content of Conjecture~\ref{conj:poly} is the claim that the
complexification or CoV \emph{always} restores ESA for polynomial
systems.  This is verified for $\dot{x}=x^{2}$ (reciprocal CoV) and
$\dot{x}=x$ (logarithmic CoV), and for the five test cases in
Section~\ref{sec:tests}.  A general proof would require showing that
every polynomial vector field admits a rational or logarithmic change of
variables rendering the flow complete, which is related to the
Painlev\'e property in the theory of integrable
systems~\cite{GORIELY2000422}.
\end{remark}

\begin{conjecture}[Higher-dimensional systems]
For coupled $M$-dimensional polynomial ODE systems, the singularity
detector (Section~\ref{sec:algorithm}) can be extended beyond 1D
quadrature by:
\begin{enumerate}
  \item Computing the blow-up time via contour integration in $\Cset^{M}$.
  \item Using the Nelson criterion on the transformed (CoV) system.
  \item Applying the BRST projection of \texttt{fock\_sirk} to identify
    the physical subspace where the extended Hamiltonian is ESA.
\end{enumerate}
\end{conjecture}

%% ====================================================================
\section*{Acknowledgments}
%% ====================================================================

The algorithmic pipeline (\texttt{ode\_sirk}) and its Lean~4
formalization (\texttt{timepiece/Singularity}) were developed as part of
the \texttt{unfer} project for quantum field theory simulations using
nested Fock spaces and rational Krylov methods.

%% ====================================================================
\appendix
\section{Build and Run Sequence}\label{app:build}
%% ====================================================================

The \texttt{ode\_sirk} crate is built and validated as follows:

\begin{enumerate}
  \item \texttt{cargo new ode\_sirk --lib} (sibling to
    \texttt{fock\_sirk}).
  \item Add dependencies: \texttt{nested\_fock\_algebra},
    \texttt{num\_complex}, \texttt{nalgebra}, \texttt{serde},
    \texttt{rustc-hash}, \texttt{thiserror},
    \texttt{quantrs2-symengine-pure}.
  \item Implement Wick's algorithm in \texttt{poly.rs} with strict unit
    tests verifying exact commutators ($1/\sqrt{2}$ scaling).
  \item Implement \texttt{hamiltonian.rs} Weyl quantization logic.
  \item Implement classical ODE solver in \texttt{flow.rs} with blow-up
    thresholds ($R_{\max}=10^{6}$, $\Delta t_{\min}=10^{-14}$).
  \item Connect \texttt{HamiltonianSpec::OdeSystem} in
    \texttt{unfer\_protocol}.
  \item Implement \texttt{unfer\_ffi} bindings:
    \texttt{uk\_ode\_analyze}, \texttt{uk\_ode\_measure\_original}.
  \item \texttt{cargo test -p ode\_sirk} to validate against Nelson's
    theoretical baselines.
\end{enumerate}

The Lean formalization is checked with:
\begin{lstlisting}[language={},caption={Lean~4 build}]
cd timepiece/Singularity
lake build
\end{lstlisting}

\begin{thebibliography}{99}

\bibitem{Sudarshan1976Interaction}
E.~C.~G.~Sudarshan, ``Interaction of classical and quantum systems,''
\textit{Pramana}, 1976.

\bibitem{PhysRevResearch.2.043102}
A.~B.~et~al., ``Koopman--von Neumann approach to quantum simulation,''
\textit{Phys.\ Rev.\ Research} \textbf{2}, 043102, 2020.

\bibitem{BEVANDA2021197}
P.~Bevanda, S.~Sosnowski, S.~Hirche, ``Koopman operator theory for
nonlinear dynamical systems,'' \textit{Annual Reviews in Control},
\textbf{52}, 197--230, 2021.

\bibitem{schmuxfcdgen2012unbounded}
K.~Schm\"udgen, \textit{Unbounded Self-Adjoint Operators on Hilbert
Space}, Springer, 2012.

\bibitem{book}
L.~Pedro, \textit{Unitary Inference}, 2024.

\bibitem{cassinelli2004theory}
G.~Cassinelli, E.~De~Vito, A.~Levrero, \textit{The Theory of Group
Representations}, Springer, 2004.

\bibitem{SIMMS1971283}
D.~J.~Simms, ``Bargmann's theorem and the Stone--von Neumann
theorem,'' \textit{Comm.\ Math.\ Phys.}, \textbf{20}, 283, 1971.

\bibitem{GORIELY2000422}
A.~Goriely, ``The Painlev\'e property for ODEs,'' in \textit{The
Painlev\'e Property}, R.~Conte, ed., Springer, 2000.

\bibitem{cmpux2f1103859517}
M.~M\"antoiu, V.~Mantoiu, ``Essential self-adjointness via auxiliary
operators,'' \textit{Comm.\ Math.\ Phys.}, 2017.

\bibitem{faris2006self}
W.~G.~Faris, \textit{Self-Adjoint Operators}, Springer, 2006.

\bibitem{hashimoto_shift-invert_2019}
K.~Hashimoto, T.~Nodera, ``Shift-invert rational Krylov method for
matrix exponentials,'' \textit{SIAM J.\ Sci.\ Comput.}, 2019.

\bibitem{hashimoto_shift-invert_2016}
K.~Hashimoto, T.~Nodera, ``Rational Krylov methods for ODEs,''
\textit{J.\ Comput.\ Appl.\ Math.}, 2016.

\bibitem{3Deuler}
T.~Hou, ``Finite-time singularity of the 3D Euler equations,'' 2024.

\bibitem{deng_hilberts_2025}
W.~Deng, ``Hilbert's sixth problem and Navier--Stokes,'' 2025.

\bibitem{article}
A.~B.~et~al., ``Koopman operator in classical mechanics,'' 2021.

\end{thebibliography}

\end{document}
